\documentclass[11pt,a4paper]{article}
\usepackage[a4paper,margin=2.2cm]{geometry}
\usepackage{fontspec}
\setmainfont{Times New Roman}
\setmonofont{Consolas}[Scale=0.86]
\usepackage{amsmath,amssymb,mathtools}
\usepackage{unicode-math}
\setmathfont{Latin Modern Math}
\usepackage{booktabs}
\usepackage{array}
\usepackage{enumitem}
\usepackage{titlesec}
\usepackage{xcolor}
\usepackage{mdframed}
\usepackage[hidelinks]{hyperref}

\definecolor{acc}{HTML}{7A2E1E}
\definecolor{dim}{HTML}{555555}
\definecolor{boxbg}{HTML}{F4F1EC}

\titleformat{\section}{\large\bfseries\color{acc}}{\thesection.}{0.5em}{}
\titlespacing{\section}{0pt}{1.1em}{0.4em}
\titleformat{\subsection}{\bfseries}{\thesubsection.}{0.5em}{}
\titlespacing{\subsection}{0pt}{0.8em}{0.25em}

\setlist[itemize]{leftmargin=1.2em,itemsep=0.15em,topsep=0.3em}
\setlist[enumerate]{leftmargin=1.4em,itemsep=0.15em,topsep=0.3em}
\renewcommand{\arraystretch}{1.15}
\setlength{\parskip}{0.45em}
\setlength{\parindent}{0pt}

\newcommand{\qq}[1]{\texttt{\small #1}}
\newcommand{\note}[1]{{\color{dim}\small #1}}

% Врезка для вывода или ключевого следствия.
\newmdenv[
  backgroundcolor=boxbg, linewidth=0pt, skipabove=0.7em, skipbelow=0.7em,
  innerleftmargin=0.9em, innerrightmargin=0.9em,
  innertopmargin=0.6em, innerbottommargin=0.6em
]{keybox}

\DeclareMathOperator*{\argmax}{arg\,max}
\DeclareMathOperator*{\argmin}{arg\,min}
\DeclareMathOperator{\KL}{KL}
\newcommand{\R}{\mathbb{R}}
\newcommand{\E}{\mathbb{E}}
\newcommand{\Var}{\operatorname{Var}}

\begin{document}

{\Large\bfseries Выравнивание и управление поведением}

\vspace{-0.5em}
\note{vlm-recipes~· подзадача 3}

\section{Чего не умеет SFT}

SFT максимизирует правдоподобие эталонного ответа. Это учит модель
\emph{воспроизводить} хорошие ответы, но ничего не говорит о том, что
один ответ лучше другого. Между тем именно это нам и нужно: правильных
ответов на «сформулируй мне гипотезу» много, и различаются они оттенком —
отдаёт ли модель решение студенту или выдаёт готовый текст.

Формально: SFT видит только положительные примеры и никогда
не наблюдает, что \emph{не} следует говорить. Градиент кросс-энтропии
поднимает вероятность эталона, но не опускает вероятность правдоподобной
альтернативы. Выравнивание по предпочтениям закрывает именно этот пробел.

\section{Классический путь: RLHF}

Сначала по парам $(y_w \succ y_l)$ обучается модель награды. Вероятность
предпочтения задаётся моделью Брэдли–Терри:
\[
p(y_w \succ y_l \mid x) = \sigma\big(r(x,y_w) - r(x,y_l)\big),
\qquad \sigma(z) = \frac{1}{1+e^{-z}} .
\]
Затем политика оптимизируется методом PPO при штрафе за отход
от исходной модели:
\[
\max_{\pi}\;
\E_{x\sim D,\,y\sim\pi}\big[r(x,y)\big]
\;-\;\beta\,\KL\!\big(\pi(\cdot\mid x)\,\|\,\pi_{\text{ref}}(\cdot\mid x)\big).
\]
Штраф необходим: без него политика уходит в вырожденные тексты, дающие
высокую награду и бессмысленные для человека.

Схема работает, но требует четырёх моделей в памяти (политика, опорная,
награды, критик) и капризна в настройке.

\section{DPO: вывод}

Ключевое наблюдение Rafailov и др. (2023): задача выше решается
аналитически, и это позволяет обойтись без модели награды вовсе.

\subsection{Шаг 1. Оптимальная политика}

Для фиксированного $x$ максимизация по $\pi$ при KL-штрафе имеет
известное решение больцмановского вида:
\[
\pi^*(y\mid x) \;=\; \frac{1}{Z(x)}\,
\pi_{\text{ref}}(y\mid x)\,\exp\!\Big(\tfrac{1}{\beta}\,r(x,y)\Big),
\qquad
Z(x) = \sum_{y}\pi_{\text{ref}}(y\mid x)\,e^{r(x,y)/\beta}.
\]

\subsection{Шаг 2. Обращение}

Логарифмируем и выражаем награду через политику:
\[
r(x,y) \;=\; \beta\,\log\frac{\pi^*(y\mid x)}{\pi_{\text{ref}}(y\mid x)}
\;+\;\beta\,\log Z(x).
\]

\begin{keybox}
Награда оказалась \emph{выражена через саму политику}. Отдельная модель
награды не нужна: любая политика неявно задаёт награду, и наоборот.
\end{keybox}

\subsection{Шаг 3. Подстановка}

Подставляем в Брэдли–Терри. Слагаемое $\beta\log Z(x)$ одинаково для
$y_w$ и $y_l$ и сокращается в разности — вместе с ним исчезает
единственная неберущаяся сумма:
\[
p(y_w \succ y_l \mid x) = \sigma\!\left(
\beta\log\frac{\pi_\theta(y_w\mid x)}{\pi_{\text{ref}}(y_w\mid x)}
-\beta\log\frac{\pi_\theta(y_l\mid x)}{\pi_{\text{ref}}(y_l\mid x)}
\right).
\]

Максимизация правдоподобия этих предпочтений и есть DPO:
\[
\mathcal{L}_{\text{DPO}} = -\,\E_{(x,y_w,y_l)}
\left[\log\sigma\!\left(
\beta\log\frac{\pi_\theta(y_w\mid x)}{\pi_{\text{ref}}(y_w\mid x)}
-\beta\log\frac{\pi_\theta(y_l\mid x)}{\pi_{\text{ref}}(y_l\mid x)}
\right)\right].
\]

\subsection{Что делает градиент}

Обозначим неявный отступ $\hat\Delta = \hat r_\theta(x,y_w) - \hat r_\theta(x,y_l)$.
Тогда
\[
\nabla_\theta \mathcal{L}_{\text{DPO}} = -\beta\,
\underbrace{\sigma(-\hat\Delta)}_{\text{вес примера}}
\Big[\nabla_\theta\log\pi_\theta(y_w\mid x)
- \nabla_\theta\log\pi_\theta(y_l\mid x)\Big].
\]

Вес $\sigma(-\hat\Delta)$ велик там, где модель \emph{ошибается}
в предпочтении, и близок к нулю там, где она уже права. Обучение само
концентрируется на трудных парах — отсюда устойчивость метода
и малое требуемое число примеров.

\section{ORPO: один этап вместо двух}

DPO требует предварительного SFT и опорной модели в памяти. ORPO
(Hong и др., 2024) обходится без обоих, добавляя к обычному SFT
слагаемое с отношением шансов.

Шансы последовательности:
\[
\text{odds}_\theta(y\mid x) = \frac{\pi_\theta(y\mid x)}{1-\pi_\theta(y\mid x)},
\qquad
\text{OR} = \frac{\text{odds}_\theta(y_w\mid x)}{\text{odds}_\theta(y_l\mid x)} .
\]
Полная функция потерь:
\[
\mathcal{L}_{\text{ORPO}} = \mathcal{L}_{\text{SFT}}
- \lambda\,\log\sigma\big(\log \text{OR}\big).
\]

Первое слагаемое тянет к эталону, второе разводит выбранный
и отвергнутый. Отношение шансов выбрано вместо отношения вероятностей
намеренно: оно наказывает мягче и не разрушает уже усвоенное на SFT.

\begin{keybox}
Практический смысл: ORPO даёт один прогон вместо двух и не держит
вторую копию модели. При ограниченных ресурсах это метод выбора,
а качество сопоставимо с SFT + DPO.
\end{keybox}

\section{SimPO: без опорной модели}

Meng и др. (2024) замечают, что неявная награда DPO не согласована
с тем, чем на самом деле оценивается генерация. Модель порождает
последовательности, максимизируя \emph{среднюю} логвероятность, тогда
как награда DPO — сумма по токенам, и потому систематически предпочитает
длинные ответы.

SimPO нормирует награду на длину и обходится без опорной модели:
\[
r_{\text{SimPO}}(x,y) = \frac{\beta}{|y|}\log\pi_\theta(y\mid x),
\qquad
\mathcal{L} = -\log\sigma\big(r(x,y_w) - r(x,y_l) - \gamma\big).
\]

Порог $\gamma > 0$ требует, чтобы выбранный превосходил отвергнутого
не просто, а с запасом. Побочный эффект — метод не раздувает длину
ответов, характерную для DPO.

\section{KTO: когда пар нет}

Собрать пары дорого: нужно два ответа на один промпт и сравнение между
ними. KTO (Ethayarajh и др., 2024) обходится одиночными примерами
с бинарной меткой «хорошо» или «плохо».

\subsection{Конструкция}

Неявная награда та же, что в DPO, — логарифм отношения политики
к опорной:
\[
r_\theta(x,y) = \log\frac{\pi_\theta(y\mid x)}{\pi_{\text{ref}}(y\mid x)}.
\]

Вторую половину пары заменяет \textbf{точка отсчёта} $z_0$ — среднее
отклонение политики от опорной, оцениваемое по батчу. Технически это
KL-дивергенция, посчитанная на перемешанных парах: ответ $y'$ берётся
от чужого промпта, чтобы оценка не зависела от конкретного примера:
\[
z_0 = \KL\big(\pi_\theta(y'\mid x)\,\|\,\pi_{\text{ref}}(y'\mid x)\big).
\]
Величина $z_0$ трактуется как константа и в градиенте не участвует.

\subsection{Функция ценности}

Ценность ответа измеряется относительно $z_0$, причём знак сравнения
зависит от метки:
\[
v(x,y) =
\begin{cases}
\lambda_D\,\sigma\big(\beta\,(r_\theta(x,y) - z_0)\big), & y \text{ помечен «хорошо»},\\[0.4em]
\lambda_U\,\sigma\big(\beta\,(z_0 - r_\theta(x,y))\big), & y \text{ помечен «плохо»}.
\end{cases}
\]
Для хорошего ответа награда должна \emph{превышать} точку отсчёта, для
плохого — \emph{отставать} от неё. Функция потерь — недобор ценности:
\[
\mathcal{L}_{\text{KTO}} = \E_{x,y\sim D}\big[\lambda_y - v(x,y)\big],
\]
где $\lambda_y$ равно $\lambda_D$ или $\lambda_U$ в зависимости от метки.

Сигмоида здесь играет роль функции ценности из теории перспектив
Канемана и Тверски: она вогнута на приобретениях и выпукла на потерях,
то есть отдача убывает по мере удаления от точки отсчёта в обе стороны.

\subsection{Почему пара не нужна}

Сопоставим с DPO. Там ответ сравнивается с \emph{конкретным}
отвергнутым:
\[
\hat r_\theta(x, y_w) - \hat r_\theta(x, y_l).
\]
В KTO ответ сравнивается со \emph{средним} поведением модели:
\[
r_\theta(x,y) - z_0 .
\]

\begin{keybox}
Точка отсчёта $z_0$ заменяет вторую половину пары. Сравнивать есть
с чем, даже когда про каждый пример известно только «хорошо» или
«плохо»: вместо «лучше ли этот ответ того» задаётся вопрос «лучше ли
этот ответ того, что модель делает в среднем».
\end{keybox}

\subsection{Асимметрия}

Веса $\lambda_D$ и $\lambda_U$ различны, и обычно $\lambda_U > \lambda_D$.
Формально это из теории перспектив: потери переживаются острее
приобретений. Практически же это рычаг под несбалансированные данные.
Жалоб на ответы всегда меньше, чем лайков, и больший вес на «плохих»
примерах компенсирует перекос — авторы рекомендуют подбирать отношение
$\lambda_D n_D / \lambda_U n_U$ в диапазоне от 1 до 4/3, где $n_D$
и $n_U$ — число хороших и плохих примеров.

Отсюда главное практическое следствие: данные можно собирать
из продакшена. Лайк или жалоба на ответ уже является меткой, и никакой
разметки пар не требуется.

\section{Сводка}

\begin{center}
\begin{tabular}{@{}llcp{5.0cm}@{}}
\toprule
\textbf{Метод} & \textbf{Данные} & \textbf{Опорная} & \textbf{Когда брать} \\
\midrule
DPO & пары & да & есть SFT-модель и память на две копии \\
ORPO & пары & нет & один этап, ограниченные ресурсы \\
SimPO & пары & нет & важна длина ответа \\
KTO & метки & да$^*$ & пары собрать дорого \\
\bottomrule
\end{tabular}
\end{center}

\note{$^*$ KTO использует опорную модель, но допускает несбалансированные данные.}

\section{Векторы управления}

Отдельный класс приёмов: изменение поведения \emph{без обучения вовсе}.

\subsection{Гипотеза линейного представления}

Эмпирически установлено, что абстрактные свойства текста — тональность,
правдивость, готовность отказать — кодируются в пространстве активаций
приблизительно линейно: существует направление $v$, сдвиг вдоль которого
усиливает свойство.

Оценка направления — разность средних по двум контрастным множествам:
\[
v_\ell \;=\; \frac{1}{|P|}\sum_{p\in P} h_\ell(p)
\;-\;\frac{1}{|N|}\sum_{n\in N} h_\ell(n),
\]
где $h_\ell(\cdot)$ — активация на слое $\ell$, усреднённая по позициям.
Применение при генерации:
\[
h_\ell \;\leftarrow\; h_\ell + s\,\frac{v_\ell}{\lVert v_\ell\rVert}.
\]

\subsection{Практика}

Нормировка делает коэффициент $s$ сопоставимым между моделями и слоями.
Слой берётся из середины сети: ранние кодируют поверхностную форму,
поздние — конкретные токены, абстрактные свойства живут между ними.
Значения $s$ выше 3--4 у большинства моделей разрушают связность.

\begin{keybox}
Главное практическое следствие: если вектор не даёт эффекта ни при каком
$s$, то и дообучение на тех же примерах, скорее всего, не поможет —
свойство не выделяется линейно, и проблема в постановке задачи. Это
самая дешёвая проверка достижимости поведения: десятки пар и один проход
вместо дней сбора данных.
\end{keybox}

\section{Порядок применения}

Приёмы упорядочены по стоимости; двигаться следует сверху вниз
и останавливаться, как только результат достигнут.

\begin{center}
\begin{tabular}{@{}llp{6.6cm}@{}}
\toprule
\textbf{Приём} & \textbf{Стоимость} & \textbf{Ограничение} \\
\midrule
системный промпт & 0 & перебивается текстом пользователя \\
префилл ответа & 0 & задаёт форму, не содержание \\
грамматика при декодировании & 0 & гарантирует формат, не факты \\
модель-судья & +1 вызов & удвоенная задержка \\
вектор управления & часы & грубый, сдвигает поведение целиком \\
SFT & дни & нужны эталонные ответы \\
DPO / ORPO & дни & нужны пары предпочтений \\
\bottomrule
\end{tabular}
\end{center}

Отдельно стоит помнить, что значительная часть правил вообще
не должна доходить до модели. Запрет на действие закрывается тем, что
инструмента не существует; проверяемый факт — детерминированной
проверкой в коде; формат — грамматикой при декодировании. Модели следует
оставлять только то, что требует понимания смысла.

\section{Префилл}

Приём, недооценённый чаще прочих. Начало ответа задаётся принудительно,
и генерация продолжает написанное:

\begin{quote}\qq{%
messages = [\{"role": "user", ...\},\\
\hspace*{2.4em}\{"role": "assistant", "content": "Уточняющий вопрос:"\}]\\
apply\_chat\_template(messages, add\_generation\_prompt=False)}
\end{quote}

Флаг выключен: служебные токены начала реплики уже поставлены шаблоном.

Механизм отличается от инструкции принципиально. Инструкция входит
в условие $x_{<t}$ наравне с текстом пользователя и конкурирует с ним.
Префилл же \emph{фиксирует} первые токены ответа, то есть меняет саму
точку, из которой разворачивается авторегрессия. Модель не решает,
задавать ли уточняющий вопрос — она уже его задаёт.

\end{document}
